% This file was converted to LaTeX by Writer2LaTeX ver. 1.9.9
% see http://writer2latex.sourceforge.net for more info
\documentclass{article}
\usepackage{calc,amsmath,amssymb,amsfonts}
\usepackage[T1]{fontenc}
\usepackage[italian]{babel}
\usepackage[style=numeric,backend=biber]{biblatex}
\author{HP}
\date{2026-07-26}
\begin{document}
\section{Prefazione}
Ci sono libri che insegnano un linguaggio di programmazione.

Altri spiegano il funzionamento di un computer.

Altri ancora raccontano la storia dell'informatica o dei suoi protagonisti.

Questo libro non nasce con nessuno di questi obiettivi.

L'idea che lo ha generato è molto semplice: mostrare come nasce davvero un videogioco.

Non un progetto costruito appositamente per spiegare alcune tecniche di programmazione, ma un videogioco reale,
sviluppato passo dopo passo, affrontando gli stessi problemi, gli stessi dubbi, gli stessi errori e gli stessi limiti
che ogni programmatore incontra quando trasforma un'idea in un programma funzionante.

Il Commodore 64 rappresenta il contesto ideale per questo percorso.

Le sue risorse limitate obbligano a riflettere su ogni scelta progettuale. Ogni byte di memoria, ogni ciclo di
elaborazione, ogni caratteristica dell'hardware contribuisce a rendere evidente qualcosa che, sui computer moderni,
spesso rimane nascosto: programmare significa soprattutto ragionare.

Per questo motivo il libro non segue la struttura di un manuale tradizionale.

Gli argomenti non vengono presentati secondo un ordine accademico e nessun capitolo nasce con l'intento di illustrare
una particolare istruzione del BASIC o una specifica tecnica di programmazione.

Ogni nuovo concetto compare soltanto quando il videogioco lo richiede.

È il progetto a determinare il percorso di apprendimento, non il contrario.

Accanto alla componente tecnica ne esiste un'altra, altrettanto importante.

La realizzazione del videogioco viene raccontata attraverso una storia vera.

Non perché la storia personale dell'autore costituisca il centro dell'opera, ma perché è proprio da quella storia che
nasce il progetto. Le persone incontrate, gli insegnamenti ricevuti, gli entusiasmi condivisi e le esperienze vissute
rappresentano il contesto umano dal quale prende forma il percorso tecnico che il lettore seguirà nelle pagine
successive.

Narrazione e programmazione procedono quindi insieme.

La prima fa emergere i problemi.

La seconda mostra come affrontarli.

L'obiettivo non è soltanto insegnare a scrivere codice, ma accompagnare il lettore nel modo di pensare che precede ogni
riga di codice.

Questo libro è anche il risultato di un lungo lavoro di progettazione editoriale.

Prima ancora di scrivere il primo capitolo è stato necessario definire con precisione l'identità dell'opera, la sua
struttura, il percorso didattico e le regole che ne avrebbero guidato la stesura. In questa fase, l'autore ha
utilizzato ChatGPT come interlocutore critico e come supporto nella revisione, nella verifica della coerenza e
nell'organizzazione delle idee. Le scelte contenute in queste pagine rimangono interamente dell'autore; l'intelligenza
artificiale ha rappresentato uno strumento di confronto, analisi e affinamento, analogo a quello che un tempo avrebbe
potuto offrire un editor durante la preparazione di un manoscritto.

Se, al termine della lettura, il lettore avrà la sensazione di aver partecipato davvero alla nascita di un videogioco e,
insieme ad essa, avrà acquisito un modo più consapevole di affrontare i problemi della programmazione, allora questo
libro avrà raggiunto il proprio scopo.

L'autore


\bigskip
\end{document}
